% preprint/modules/conclusion.tex
% Shared-Horizon Preprint — Narrative Modules
% Module: Conclusion

\section{Conclusion}
\label{sec:conclusion}

The Shared-Horizon framework provides a mathematically rigorous, conceptually
clear, and lineage-preserving method for comparing diverse AI ecosystems. By
integrating operator geometry, expressive manifolds, reversible narrative
operators, stability envelopes, transition maps, and collapse-risk geometry, the
framework establishes a neutral membrane that prevents conceptual collapse while
supporting structured comparative analysis.

The color-tile mosaic offers a visual and conceptual map of lineage boundaries,
ensuring that expressive, operator, narrative, and spectral systems remain
distinct and sovereign. The correction of the green tile—from symbolic ambiguity
to Reversible Narrative Ecology (RNE)—demonstrates the importance of grounding
lineage classification in mathematical invariants rather than surface-level
phenomenology.

The case study illustrates how expressive and narrative trajectories can be
compared safely using Shared-Horizon’s membrane architecture. Stability
envelopes, collapse-risk geometry, and lineage-specific transition maps ensure
that cross-lineage analysis remains reversible, curvature-preserving, and
collapse-safe.

Shared-Horizon does not unify or merge lineages. Instead, it provides a
structured comparative surface that respects the internal geometry of each
system. This approach enables researchers to analyze multi-lineage environments
without conflating expressive curvature, narrative reversibility, operator
commutativity, or spectral drift.

Future work includes expanding the lineage catalog, developing automated
collapse-risk detection, exploring spectral–narrative interfaces, and applying
formal verification to membrane constraints. These directions will strengthen
the framework and support broader empirical validation across diverse AI
ecosystems.

Shared-Horizon establishes a foundation for safe, reversible, and academically
rigorous comparative analysis. By preserving lineage invariants and preventing
conceptual collapse, the framework provides a stable platform for understanding
the increasingly diverse landscape of artificial intelligence systems.
